\documentclass[11pt,a4paper,twoside]{article}
\usepackage[utf8]{inputenc}
\usepackage[english]{babel}
\usepackage{actuarialangle, actuarialsymbol}

\begin{document}
Define exponential family and its components

A distribution belongs to the exponential family if its PDF is of the form
\(
f(y) = c(y, \phi) \exp \left\{ \frac{y\theta - a(\theta)}{\phi} \right\},
\)
where \(\theta\) and \(\phi\) are parameters. 
\(\theta\): canonical parameter 
\(\phi\): dispersion parameter. 
\( c(y, \phi) \): roughly a normalization factor

Their moments are given via
\( E(y) = \dot{a}(\theta), \quad \text{Var}(y) = \phi \ddot{a}(\theta) \).
One defines \( V(\mu) = \ddot{a}(\theta) \) the variance function.

Some define dispersion as \( \frac{\phi}{w} \) with weight \( w \), this is typical in a context of frequency.

Examples:

INSERT IMAGE FROM slide 40 


Describe the Method of Moments for parameter estimation

The Method of Moments (MoM) estimates parameters \(\theta\) of a probability distribution by equating theoretical moments to sample moments.  
General Procedure:
1. Compute the theoretical moments \( m_k \) in terms of parameters \( \theta \).  
\(
m_k = E[X^k]
\)
2. Compute the sample moments \( \hat{m}_k \) from observed data \( X_1, X_2, ..., X_n \).  
\(
\hat{m}_k = \frac{1}{n} \sum_{i=1}^{n} X_i^k
\)
3. Solve the system of equations: 
\(m_k(\theta_1, \theta_2, \dots, \theta_r) = \hat{m}_k, \quad k = 1, 2, \dots, r\) 
to obtain parameter estimates \( \hat{\theta} \).
## **Examples**  

Example Normal Distribution:
\( X \sim N(\mu, \sigma^2) \)  
Theoretical moments:  
\(
E[X] = \mu, \quad E[X^2] = \mu^2 + \sigma^2
\)
Sample moments:  
\(
\hat{m}_1 = \frac{1}{n} \sum X_i, \quad \hat{m}_2 = \frac{1}{n} \sum X_i^2
\)
Solve for parameters:  
\(
\hat{\mu} = \hat{m}_1
\)
\(
\hat{\sigma}^2 = \hat{m}_2 - \hat{m}_1^2
\)
This simplifies in the GLM via \( \dot{a}(\theta) = \bar{\mu}, \phi\ddot{a}(\theta) = \hat{\sigma}^2 \).

Key Properties
Consistency: As \( n \to \infty \), \( \hat{\theta} \to \theta \) if moments exist and are unique.  
Bias: MoM estimators may be biased.  
Efficiency: MoM is generally less efficient than MLE but simpler.  


Describe MLE for exponential family distributions

Under the i.i.d. assumption, the probability of observing the vector \( y \) of outcomes is:
\( f(y; \theta, \phi) = \prod_{i=1}^n f(y_i;\theta,\phi) \)
which is usually easier to maximize via
\( l(\theta, \phi) = \sum^n_{i=1} \ln f(y_i;\theta,\phi) \)
In particular for exp. families

\(
\sum_{i=1}^{n} \left\{ \ln c(y_i, \phi) + \frac{y_i \theta - a(\theta)}{\phi} \right\} = \frac{n \left\{ \bar{y} \theta - a(\theta) \right\}}{\phi} + \sum_{i=1}^{n} \ln c(y_i, \phi).
\)

Differentiating \(\ell(\theta, \phi)\) with respect to \(\theta\) and equating to zero leads to the first order condition for 
likelihood maximization
\(
\frac{n \{\bar{y} - \dot{a}(\theta)\}}{\phi} = 0 \quad \Rightarrow \quad \dot{a}(\theta) = \bar{y}.
\)

Hence the MLE of \(\theta\) is obtained by finding \(\theta\) such that \(\dot{a}(\theta) \equiv \mu\) equals the sample mean \(\bar{y}\). 
Thus for any exponential family distribution, \(\hat{\mu} = \bar{y}\).


What are characteristics of the MLE method?

Invariance: If \( h \) is a monotonic function and \( \hat{\theta} \) is the MLE of \( \theta \), then \( h(\hat{\theta}) \)
is the MLE of \( h(\theta) \).

Asymptotically unbiased: The expected value \( \mathbb{E}[\hat{\theta}] \) approaches \( \theta \) as the
sample size increases.

Consistent: As the sample size increases the probability distribution \( \hat{\theta} \) collapses on \( \theta \).
Minimum variance: In the class of all estimators, for large samples, \( \hat{\theta} \) has the minimum variance and 
is therefore the most precise estimate possible.

Small-sample bias. For small \( n, \hat{\theta} \) is often biased. However for large sample sizes, this is neglible.
Computationally difficult. MLEs are not always in closed form, in which case they need to be computed iteratively.
However on modern equipment this is rarely an issue.


State the classical linear model (CLM) in vector notation

In order to determine the unknown parameters \(\beta\), we collect data
\[ y_i, \; \mathbf{x}_i = (1, x_{i1}, \ldots, x_{ik})^\top, \]
with error observations according to
\[ y_i = \mathbf{x}_i^\top \boldsymbol{\beta} + \epsilon_i, \quad i = 1, 2, \ldots, n. \]
Vector notation is appropriate:

\[ 
\mathbf{y} = \begin{pmatrix}
y_1 \\
\vdots \\
y_n
\end{pmatrix}, \quad \boldsymbol{\epsilon} = \begin{pmatrix}
\epsilon_1 \\
\vdots \\
\epsilon_n
\end{pmatrix}, \quad \mathbf{X} = \begin{pmatrix}
\mathbf{x}_1^\top \\
\vdots \\
\mathbf{x}_n^\top
\end{pmatrix} = \begin{pmatrix}
1 & x_{11} & \cdots & x_{1k} \\
\vdots & \vdots & \ddots & \vdots \\
1 & x_{n1} & \cdots & x_{nk}
\end{pmatrix}.
\]

\(\mathbf{X}\) is the design matrix. One assumes that \(\text{rk} \mathbf{X} = k + 1 = p \leq n\).


The model
\( y = X \boldsymbol{\beta} + \epsilon \)
\( y \mid X \sim \mathcal{N}(X\beta, \sigma^2 I)  \)
is called the classical linear regression model, if the following assumptions hold:

\( \epsilon \sim \mathcal{N}(0, \sigma^2 I) \)
The correlation structure is also referred to as homoscedastic
while different variances for different observiations (so \( \text{Var}[\epsilon]_{i,j} \neq 0,\ i \neq j \)) would be heteroscedastic.
One says that these assumptions about \( \epsilon \) are conditionally on \( X \).
Further the errors are then uncorrelated.
\[ \text{rk}(X) = k + 1 = p \leq n \]


Derive the simple least squares estimator

This is a method for assigning values to \(\beta_0\) and \(\beta_1\) in \( y \approx \beta_0 + \beta_1x \), 
given a sample of \(n\) pairs of measurements or observations on \(y\) and \(x\) denoted \((y_i, x_i), i = 1, \ldots, n.\)
In particular, least squares determines \(\beta_0\) and \(\beta_1\) by minimizing the error sum of squares criterion:
\(
\sum_{i=1}^{n} \{y_i - (\beta_0 + \beta_1 x_i)\}^2.
\)
Using calculus, the values of \(\beta_0\) and \(\beta_1\) that minimize (4.2) are 
\(
\hat{\beta}_1 = \frac{\sum_{i=1}^{n} x_i y_i - n x \bar{y}}{\sum_{i=1}^{n} x_i^2 - n x^2}, \quad \hat{\beta}_0 = \bar{y} - \hat{\beta}_1 \bar{x},
\)
where \(\bar{y} = \sum_{i=1}^{n} y_i / n\) and \(\bar{x} = \sum_{i=1}^{n} x_i / n\). 
Predicted values based on the fit are \(\hat{y}_i = \hat{\beta}_0 + \hat{\beta}_1 x_i,\) and goodness of fit is typically measured using
\(
R^2 \equiv 1 - \frac{\sum_{i=1}^{n} (y_i - \hat{y}_i)^2}{\sum_{i=1}^{n} (y_i - \bar{y})^2}.
\)

It can be shown \(0 \leq R^2 \leq 1\) with \(R^2 = 1\) indicating a perfect linear fit. 
The \(R^2\) statistic is much abused and misinterpreted.


Derive the optimal estimator for the classical linear model (CLM) and relate it to the log-likelihood.


Minimization of \(LS\) can be achieved by inspecting the first-order conditions: 
\(\frac{\partial LS(\beta)}{\partial \beta} \bigg|_{\beta = \hat{\beta}} = -2\mathbf{X}^\top \mathbf{y} + 2 \mathbf{X}^\top \mathbf{X} \hat{\beta} \stackrel{!}{=} 0.\)

This is equivalent to the normal equations: 
\(\mathbf{X}^\top \mathbf{X} \hat{\beta} = \mathbf{X}^\top \mathbf{y}.\)

Due to the condition of full rank, the normal equations possess a unique solution, the least squares estimator, 
\(\hat{\beta} = (\mathbf{X}^\top \mathbf{X})^{-1} \mathbf{X}^\top \mathbf{y}.\)

Under the condition of normally distributed errors, i.e., \(\varepsilon \sim N(0, \sigma^2 I)\), 
the maximum likelihood estimator (MLE) equals the LSE.
This can be verified by minimizing the likelihood \(L(\beta, \sigma^2) = \frac{1}{(2\pi\sigma^2)^{n/2}} \exp \left( -\frac{1}{2\sigma^2} (y - \mathbf{X}\beta)^\top (y - \mathbf{X}\beta) \right),\)

or, equivalently, the log-likelihood

\(l(\beta, \sigma^2) = (\text{terms independent of } \beta) - \frac{1}{2\sigma^2} (y - \mathbf{X}\beta)^\top (y - \mathbf{X}\beta).\)


What is the optimal estimator in weighted least squares? 

The precision \(w_i\) of case \(i\) is signalled through \(\text{Var}(\epsilon_i) = \sigma^2 / w_i\). 
Thus \(w_i \to 0\) indicates no precision and \(w_i \to \infty\) indicates complete precision. 
Multiplying case \(i\) by \(\sqrt{w_i}\) yields 
\(y_i^* = \beta_0 \sqrt{w_i} + \beta_1 x_{i1}^* + \cdots + \beta_p x_{ip}^* + \epsilon_i^*\)
where
\(y_i^* \equiv \sqrt{w_i} \, y_i, \quad x_{ij}^* \equiv \sqrt{w_i} \, x_{ij}, \quad \epsilon_i^* \equiv \sqrt{w_i} \, \epsilon_i.\)
Thus the weighted observations are homoskedastic since \(\text{Var}(\epsilon_i^*) = \sigma^2\). Regressing \(y_i^*\) on the \(x_{ij}^*\) yields 
\(\hat{\beta}^* = (X^{*'} X^*)^{-1} X^{*'} y^* = (X'WX)^{-1}X'Wy,\)

where \(W\) is the diagonal matrix with diagonal entries \(w_i\). The estimator \(\hat{\beta}^*\) is called the weighted least squares estimator. 

Since \(\hat{\beta}^*\) is the least squares estimator computed relative to the model \(y^* = X^*\beta + \epsilon^*\) it follows that
\(\hat{\beta}^* \sim \mathcal{N} \left\{ \beta, \sigma^2 (X'WX)^{-1} \right\}.\)



State the properties of the optimal estimator for the CLM

The following properties hold if the model is correctly specified, but without any specific distributional assumptions (regarding \( \epsilon \):
\(\mathbb{E}(\hat{\beta}) = \beta\) (unbiased)
\(
\text{Cov}(\hat{\beta}) = \sigma^2 (\mathbf{X}^\top \mathbf{X})^{-1} \quad \text{with estimator} \quad \widehat{\text{Cov}}(\hat{\beta}) = \hat{\sigma}^2 (\mathbf{X}^\top \mathbf{X})^{-1}, \quad \hat{\sigma}^2 = \frac{1}{n-p} \hat{\varepsilon}^\top \hat{\varepsilon}
\)
This implies that for large \( n \) the estimator \( \hat{\beta}_n \) possesses approximately the normal distribution
\(
  N\left( \beta, \hat{\sigma}^2_n (\mathbf{X}_n^T \mathbf{X}_n)^{-1} \right).
\)
For Gaussian errors as in the CLM, the distribution of weighted distance:
\(
\frac{(\hat{\beta} - \beta)^\top (\mathbf{X}^\top \mathbf{X})(\hat{\beta} - \beta)}{\sigma^2} \sim \chi^2_p
\)
Gauss-Markov-Theorem: The LSE has the smallest variance among all linear and unbiased estimators of \(\beta\).



TODO: Check if Model selection was a topic


\end{document}
